\documentclass[11pt,openany,oneside]{article}

\usepackage{makeidx}
\usepackage{latexsym}

%estes dois packages sao para tornar automaticos a acentuacao grafica em portugues: ã, é, ô, etc...
%infelizmente e misteriosamente, isso nao funciona no meu mac, no windows funcionava...
%tendo que usar \~a para ã, etc... 

\usepackage[latin1]{inputenc}
\usepackage[portuguese]{babel}
\usepackage[T1]{fontenc}
\usepackage{amsmath}
\usepackage{amsfonts}

% essa linha foi utilizada para corrigir o erro
% auto expansion is only possible with scalable
\usepackage{lmodern}

\usepackage{titlesec}

% possibilida criar figuras 1a) e 1b), por exemplo
\usepackage{subfigure}
\usepackage[pdftex]{graphicx}

%possibilita criar figuras com legenda na lateral
\usepackage{sidecap}

\usepackage[nottoc,notlot,notlof]{tocbibind}
\usepackage{capt-of}
\usepackage[left=4.0cm, right=2.0cm, top=4.0cm,bottom=3.0cm]{geometry}%normas icontec
\usepackage{color}
\usepackage{calc}
\usepackage{titletoc}

%para criar o espacamento usual no primeiro paragrafo
\usepackage{indentfirst}

\usepackage{pifont}

% para criar os links no texto quando cita equacoes, figuras, referencias, secoes, etc...
\usepackage{hyperref}

% para colorir os links criados anteriormente
\hypersetup{ colorlinks,
linkcolor=blue,
filecolor=darkgreen,
urlcolor=blue,
citecolor=blue }
\usepackage{microtype}
\usepackage[nottoc]{tocbibind} 
\usepackage{graphicx}
\usepackage{graphicx,color}

%colocar texto do lado de figuras
\usepackage{wrapfig}


\usepackage{epsfig}
\usepackage{amssymb}
\usepackage{amsthm}
\usepackage{verbatim}

%editar captions
%\usepackage{floatrow}
%\usepackage{caption}


% possibilita inserir paginas avulsas de arquivos pdf no meio do texto
\usepackage{pdfpages}

% altera a fonte nas legendas das figuras
\usepackage[font=small,format=plain,labelfont=bf,up,textfont=it,up]{caption}



% este define o espacamento entre as linhas
% para referencia, 1.5 'e um espacamento grande, utilizado na minha tese de doutorado
\renewcommand*{\baselinestretch}{1.2}

%comando para colocar vetor linha
%deve ser usado com o comando \vec{p}+\pvec{p}'=\pvec{p}''
\newcommand{\pvec}[1]{\vec{#1}\mkern2mu\vphantom{#1}}


%\nopagenumbers

%%%%%%%%%%%%%%%
%\usepackage{fancyhdr}
%\pagestyle{fancy}

%\usepackage{xcolor}

%\setlength\headheight{26pt} %% just to make warning go away. Adjust the value after looking into the warning.
%\rhead{{\color{blue}\rule{1cm}{1cm}}}
%\rhead{\includegraphics[width=1 in]{logo_ufg1.jpg}}
%%%%%%%%%%%%%%%%%%%%%%%%%





\begin{document}
\hsize = 7.2 in
\hoffset = -0.9 in
\voffset = -0.8 in
\vsize = 9.3 in

\pagenumbering{gobble} % remove o numero da pagina

\begin{wrapfigure}{r}{0.2\textwidth}
  \vspace{-60pt}
  \begin{center}
    \includegraphics[width=1 in]{logo_ufg1.jpg}
  \end{center}
  \vspace{-20pt}
  %\caption{A gull}
  \vspace{-10pt}
\end{wrapfigure}

\subsection*{Jata�, 03/10/14. Universidade Federal de Goi�s - Campus Jata�.}

\subsection*{Curso: F�sica. Eletromagnetismo. Prof. Paulo Freitas Gomes.}

\vspace{0.2 in}
\noindent
Aluno: \_\_\_\_\_\_\_\_\_\_\_\_\_\_\_\_\_\_\_\_\_\_\_\_\_\_\_\_\_\_\_\_\_\_\_\_\_\_\_\_\_ Matr�cula:\_\_\_\_\_\_\_\_\_\_

\vspace{0.2 in}

\noindent
1) O potencial na superf�cie de uma esfera de raio $R$ � dado por:
\begin{equation}
V_0 (\theta) = k \cos 3 \theta = \dfrac{k}{5} \left[ 8P_3 (\cos \theta ) - 3 P_1 (\cos \theta ) \right] 
\nonumber
\end{equation}
onde $k$ � uma constante e $P_1 (\cos \theta ) = \cos \theta$ e $P_3 (\cos \theta ) = \frac{1}{2} ( 5\cos^3 \theta - 3\cos \theta )$ s�o os polin�mios de Legendre de ordem 1 e 3. a) Mostre que o potencial em todo o espa�o �: 
\begin{equation}
V (r,\theta) = \left\{
\begin{array}{rl}
 -\dfrac{3k}{5R} r P_1 (\cos \theta ) + \dfrac{8k}{5R^3} r^3 P_3 (\cos \theta ) & \text{se  } r \leq R  \\
\\
 -\dfrac{3kR^2}{5r^2}  P_1 (\cos \theta ) + \dfrac{8kR^4}{5r^4} P_3 (\cos \theta )  & \text{se  } r \geq R
\end{array} 
\right. \nonumber
\end{equation}
b) Assumindo que n�o h� carga nem dentro nem fora da esfera, mostre que a densidade de carga na superf�cie �:
\begin{equation}
\sigma (\theta) = \dfrac{\varepsilon_0 k}{5R} \left[ 56 P_3 (\cos \theta )-9 P_1 (\cos \theta ) \right] \nonumber
\end{equation}

%\vspace{0.1 in}
\noindent
2) Uma carga pontual $q$ de massa $m$ � solta do repouso a uma dist�ncia $d$ de um plano infinito aterrado. Mostre que o tempo necess�rio para essa carga atingir o plano �:
\begin{equation}
T = \dfrac{\pi d}{q}\sqrt{2\pi \varepsilon_0 md} \nonumber
\end{equation}
\textbf{F�rmulas para consulta}

  \begin{equation}
 \vec{p} = \int {\vec{r}\,}' \rho({\vec{r}\,}') d\tau' \quad \quad V_{dip}({\vec{r}\,}') =\dfrac{\vec{p} \cdot \hat{r}}{4\pi \varepsilon_0r^2}    \quad \quad \vec{E}_{dip} (r,\theta) = \dfrac{p}{4\pi \varepsilon_0r^3}(2\hat{r}\cos \theta + \hat{\theta} \sin \theta)   \nonumber 
\end{equation} 

 \begin{equation}
 \int_a^b x^n dx = \dfrac{b^{n+1}-a^{n+1}}{n+1}  \quad \quad \vec{F} = -\dfrac{1}{4\pi \varepsilon_0} \dfrac{q^2}{(2z)^2}\hat{z} 
 \quad \quad 
 V (r,\theta) = \left\{
\begin{array}{rl}
 V_d (r,\theta) = \sum\limits_{n=0}^{\infty} A_n r^n P_n (\cos \theta) & \text{se  } r \leq R  \\
\\
V_f (r,\theta) =\sum\limits_{n=0}^{\infty} \dfrac{B_n}{r^{n+1}} P_n (\cos \theta)  & \text{se  } r \geq R
\end{array} 
\right.
 \nonumber 
\end{equation} 
%\vspace{0.01 in} 
 
\begin{equation}
  A_n = \dfrac{2n+1}{2R^n} \int_0^{\pi} V_0(\theta) P_n (\cos \theta) \sin \theta d\theta  \quad \quad V_d(R,\theta)= V_f(R,\theta) \quad \quad \sigma (\theta) = \varepsilon_0 \sum_{n=0}^{\infty} (2n+1) A_n R^{n-1} P_n (\cos \theta )   \nonumber 
\end{equation} 



\begin{equation}
 \int_0^{\pi} P_n  (\cos \theta ) P_m  (\cos \theta ) \sin \theta d\theta  = \left\{
\begin{array}{rl}
 0& \text{se  } n \neq m  \\
\\
 \dfrac{2}{2n+1}& \text{se  } n = m
\end{array} 
\right.
\quad \quad
\int_d^0 \sqrt{\dfrac{x}{d-x}}dx = -d\dfrac{\pi}{2} \quad \quad \sum_i \vec{F}_i = m\dfrac{d^2\vec{r}}{dt^2} 
 \nonumber 
\end{equation}

\begin{equation}
 \left(  \dfrac{\partial V_d }{\partial r}- \dfrac{\partial V_f }{\partial r} \right)\biggr\rvert_{r=R}   = -\dfrac{\sigma(\theta)}{\varepsilon_0}  \quad \quad \nabla^2 V(\vec{r}) = \rho(\vec{r}) \quad \quad V(\vec{r}) = \dfrac{1}{4\pi \varepsilon_0} \sum_{n=0}^{\infty} \dfrac{1}{r^{n+1}} \int (r')^n P_n (\cos \theta) \rho({\vec{r}\,}')d\tau'  \nonumber 
\end{equation} 






 
\end{document}



